
\section{Introducción} \label{sec: Intro}
	
	El proyecto se desarrolla en el ámbito académico de la Universidad Siglo 21. Nacido como una actividad trasversal de Abogacía y Sistemas, para dar una solución a la gestión de los reclamos relacionados con la defensa del consumidor por parte de asociaciones y ONGs vinculadas a los derechos del consumidor de la Provincia de Córdoba, pudiendo ser usada por otras organizaciones del resto del país. Que permitirá agilizar los tiempos de gestión y dar visibilidad a los trámites asociados a estas reclamaciones.\\
	
	La asociación con la que se trabajará se llama \textbf{Novo Derecho SAS}.
	
	\subsection{Antecedentes} \label{subsec: Antec}
		
		``El régimen normativo argentino de defensa de los usuarios y consumidores fue
		conformándose paulatinamente a través de sucesivas etapas, propias de todo "proceso". La primera de ellas se verificó en el mes de octubre de 1993, oportunidad en la cual el Congreso de la Nación sancionó la ley 24240'' (\cite{Hernandez2003})
			
		Desde la creación de la ley existen en gran parte de las jurisdicciones del país, un organismo gubernamental que se encarga de la recepción y administración de los reclamos que se generan por el incumplimiento de dicha ley, pero en muchas ocasiones no da abasto para atender todas las reclamaciones. Entonces se crearon en el país, cientos de asociaciones civiles que buscan dar una solución extrajudicial a los ciudadanos que tengan problemas con las compras de bienes y/o contrataciones de servicios. Estas organizaciones están amparadas por la Ley de Defensa al Consumidor (LDC\footnote{Ley de Defensa del Consumidor. Ley N° 24.240}).\\
		
		En el caso particular de la Provincia de Córdoba existe la web Ciudadano Digital (\url{https://cidi.cba.gov.ar}) página del Gobierno de la provincia que agrupa tramites varios, entre ellos los reclamos de defensa del consumidor. Antes del inicio de la pandemia de COVID-19 las reclamaciones se hacían en la oficina de defensa al consumidor, pero por razones sanitarias esos reclamos fueron pasados a la web. Si bien el sitio es moderno, no es una plataforma destinada exclusivamente para la carga de este tipo de reclamos, teniendo el usuario que sortear impedimento que van desde, la localización de los enlaces que llevan a los formularios, hasta la terminología utilizada, demasiado apegada a al derecho, que pocos ciudadanos comprenden. 
		
		%En antecedentes se hace una reseña histórica, descripción en el tiempo de la problemática y/o los %avances o trabajos previos que se hayan realizado sobre el tema.
		
	\subsection{Descripción del área problemática} \label{subsec: DescAreaProblem}
		La web actual que se utiliza para gestionar los reclamos relacionados con la LDC en la Provincia de Córdoba no está adaptada para todos los ciudadanos, utilizando terminología muy apegada al derecho, que no todos conocen. El acceso al sitio donde se cargan los reclamos requiere cierta experiencia en el manejo de la web, y si bien en el año 2021 uno tiende a pensar que ya la mayor cantidad de personas están familiarizadas con estas tecnologías, hay una gran parte de la población que no las maneja.\\
		
		El seguimiento del reclamo ya sea por parte de los abogados o los mismos consumidores que deciden ellos mismos llevar adelante sus casos, es virtualmente imposible. El sistema sólo otorga un número de caso, pero la documentación que el usuario registra en la página no se aloja en ningún sitio “a la vista” del reclamante. No hay un comprobante de realización del trámite o algo que indique al usuario que el reclamo se cargó correctamente.\\
		
		El sistema actual no emite ninguna notificación al usuario, ante falta de información, documentación o avance del caso, provocando en ocasiones que los reclamos no prosperen, con claro perjuicio al consumidor.
		
		%Presenta en términos generales el área o situación problemática. Debe incluir aspectos como: áreas %afectadas, factores involucrados, magnitud del problema, frecuencia de presentación, tendencias %observadas, personas e instituciones involucradas.
		
	\subsection{Formulación de la problemática} \label{subsec: FormProblema}
		Teniendo en cuenta lo expuesto en el punto anterior y la realidad actual sobre los problemas sanitarios que obliga a las personas a gestionar los aspectos administrativos vía web, es que surge la pregunta: ¿Es factible crear un sitio que de una solución alternativa a la vigente? Agilizar los procesos de carga de información, ser amigable con el consumidor/usuario, sea éste experimentado o sin ningún tipo de conocimiento en materia de reclamos legales, dar visibilidad a la gestión administrativa del proceso, conocer las instancias de su reclamo.
		
		%Sin ser una regla, una forma clara de formular el problema es %dándole estructura de pregunta, tan breve como sea posible. Esta %formulación ayuda al autor a plantear los objetivos.
	